\begin{textsample}{POS Dim 4 – pr – Score 11.00 – pr/pr\_ef\_6.txt}  \label{ex:f4_pos_185}
Ensino \textbf{religioso} O \textbf{Estado} do Paraná tem sido referência para todo o Brasil pelo trabalho desenvolvido em prol da disciplina de Ensino \textbf{Religioso}.
Com o intuito de contemplar o disposto no Art. 33 da Lei de Diretrizes e Bases da Educação Nacional - LDB/96, o qual determina que a disciplina deve fomentar “ o respeito à diversidade cultural \textbf{religiosa} do Brasil vedadas quaisquer formas de proselitismo ”, é imprescindível uma imparcialidade ideológica dos professores, não direcionando os estudantes a uma determinada corrente de pensamento, seja ela \textbf{religiosa} ou não.
A disciplina de Ensino \textbf{Religioso} está presente nos currículos escolares no Brasil, assumindo diferentes formatos de acordo com os períodos históricos e a legislação vigente.
A primeira forma de inclusão dos temas \textbf{religiosos} na educação brasileira, que se perpetuou até a Constituição da República em 1891, pode ser identificada nas atividades de evangelização promovidas pela Companhia de Jesus, de confissão católica, conforme o documento nominado de Ratio Studiorum.
Com o advento da República e do ideal positivista de separação entre Estado e Igreja, todas as instituições e assuntos de ordem pública buscaram se reestruturar de acordo com o critério de laicidade interpretada no sentido de neutralidade \textbf{religiosa}.
Em 1934, a disciplina de Ensino \textbf{Religioso} passa a ser contemplada nos currículos da educação pública, salvaguardando o direito individual de liberdade de credo.
Dessa forma, o artigo da Constituição da Era Vargas que tratava do Ensino \textbf{Religioso} trazia a seguinte redação : > “ O ensino \textbf{religioso} será de frequência facultativa e ministrado de acordo com os princípios da confissão \textbf{religiosa} do aluno manifestada pelos pais ou responsáveis e constituirá matéria dos horários nas escolas públicas primárias, secundárias, profissionais e normais ” ( BRASIL,1934, art. 153 ).
Dessa forma, a Constituição de 1934, assim como as que vieram na sequência, pretendiam responder à questão da laicidade do \textbf{Estado} com o acréscimo e manutenção do caráter facultativo da disciplina, uma vez que, legalmente garantido o direito de não participar do Ensino \textbf{Religioso}, a liberdade de credo do cidadão estaria igualmente garantida.
A concepção \textbf{religiosa} desse período era, portanto, restritiva e abordava unicamente a doutrina cristã.
Somente na Constituição de 1988 em seu Art. 210 - §.1º, o teor do texto ficou mais sucinto no que diz respeito a laicidade quando afirma : “ O Ensino \textbf{Religioso}, de matrícula facultativa, constituirá disciplina dos horários normais das escolas públicas de ensino fundamental ”.
Apesar do que acontecia no Brasil até a década de 1980, mundialmente os impulsos contrários à perspectiva confessional de ensino se tornavam cada vez mais fortes.
A Declaração Universal dos Direitos Humanos, promulgada em 1948, afirmava em seu 18º artigo o seguinte : “ Toda pessoa tem o direito à liberdade de pensamento, consciência e religião ; este direito inclui a liberdade de mudar de religião ou crença e a liberdade de manifestar essa religião ou crença pelo ensino, pela prática, pelo culto e pela observância isolada ou coletivamente, em público ou em particular ”.
A possibilidade de um Ensino \textbf{Religioso} aconfessional, coerente com um \textbf{Estado} Laico[24 ] só se concretizou legalmente na redação da Lei de Diretrizes e Bases da Educação Nacional de 1996 e sua respectiva correção, em 1997, pela Lei 9.475/97.
De acordo com o artigo 33 da LDBEN, o Ensino \textbf{Religioso} recebeu a seguinte caracterização : > Art. 33 – O Ensino \textbf{Religioso}, de matrícula facultativa, é parte integrante da formação básica e constitui disciplina dos horários normais das escolas públicas de Educação Básica assegurado o respeito à diversidade \textbf{religiosa} do Brasil, vedadas quaisquer formas de proselitismo. > § 1º – Os sistemas de ensino regulamentarão os procedimentos para a definição dos conteúdos do Ensino \textbf{Religioso} e estabelecerão as normas para a habilitação e admissão de professores. > § 2º – Os sistemas de ensino ouvirão entidade civil, constituída pelas diferentes denominações \textbf{religiosas}, para a definição dos conteúdos do ensino \textbf{religioso}.
Para viabilizar a proposta de Ensino \textbf{Religioso} no Paraná, a Associação Inter \textbf{Religiosa} de Curitiba ( Assintec ), formada por um grupo de representantes das diversas organizações \textbf{religiosas} que formam a sociedade civil organizada, atua desde 1973 em conjunto com Estados e Municípios na elaboração de material pedagógico e cursos de formação continuada.
Nesse sentido, considerando o processo histórico vivenciado pelo \textbf{Estado} do Paraná, a construção dos documentos orientadores estaduais para a Educação Básica, as Diretrizes Curriculares Nacionais e a homologação da Base Nacional Comum Curricular - BNCC para o Ensino Fundamental, que define as Competências Gerais e Específicas para a Área de Ensino \textbf{Religioso}, é que se elabora este Referencial Curricular do Paraná : princípios, direitos e orientações.
É importante destacar que o documento em questão foi desenvolvido pelos técnicos pedagógicos da equipe de Currículo da Secretaria de Estado da Educação do Paraná ( SEED ), em um trabalho conjunto com a equipe pedagógica da Associação Inter \textbf{Religiosa} de Educação e Cultura ( ASSINTEC ) e com a equipe pedagógica da Secretaria Municipal de Curitiba ( SME ), representando a União Nacional dos Dirigentes Municipais de Educação ( UNDIME ).
É importante salientar que o objeto de estudo do Componente Curricular Ensino \textbf{Religioso} tem variado ao longo de sua história.
Contudo, no atual contexto da rede pública estadual, O Sagrado está definido como objeto de estudo, dessa forma possibilita o estudo da manifestação da diversidade \textbf{religiosa} e cultural concebido como a forma da religiosidade se manifestar e poder ser estudada.
Na BNCC foi adotado o conceito de Conhecimento \textbf{Religioso} como objeto de estudo da área de Ensino \textbf{Religioso}, o qual é produzido no âmbito das diferentes áreas do conhecimento científico das Ciências Humanas e Sociais, principalmente nas Ciência(s) da(s) Religião(ões), visto que essas Ciências investigam e analisam as manifestações dos fenômenos \textbf{religiosos} em diferentes culturas e sociedades.
Entende -se como manifestações do fenômeno \textbf{religioso} : as cosmovisões, linguagens, saberes, crenças, temporalidade sagrada, festas \textbf{religiosas}, mitologias, narrativas, textos, símbolos, ritos, doutrinas, tradições/organizações, práticas e princípios éticos e morais.
Os fenômenos \textbf{religiosos} em suas múltiplas manifestações são parte integrante do substrato cultural da humanidade ( BRASIL, 2017, pg. 434 ).
O desenvolvimento e a organização do Referencial Curricular do Paraná foram elaborados em consonância com as Competências Gerais da BNCC.
Para tanto, o Ensino \textbf{Religioso} deve atender os seguintes objetivos : a.
Proporcionar a aprendizagem dos conhecimentos \textbf{religiosos}, culturais e estéticos, a partir das manifestações \textbf{religiosas} percebidas na realidade dos educandos sempre contemplando as 4 matrizes \textbf{religiosas} que forma a religiosidade brasileira ( \textbf{Indígena}, Afro, Ocidental e Oriental ) ; b.
Propiciar conhecimentos sobre o direito à liberdade de consciência e de crença tanto individuais e coletivas, com o propósito de promover o conhecimento e a efetivação do que está prescrito na Declaração Universal dos Direitos Humanos ; c.
Desenvolver competências e habilidades que contribuam para o diálogo entre perspectivas \textbf{religiosas} e seculares diferentes de vida, exercitando o respeito à liberdade de concepções e o pluralismo de ideias, de acordo com a Constituição Federal ; d.
Contribuir para que os educandos construam seus sentidos pessoais de vida a partir de valores, princípios éticos e da cidadania. ( BRASIL, 2017, pg. 434 ).
Nesse sentido, as Competências Específicas apontadas para o Ensino \textbf{Religioso} na BNCC e, por consequência, presentes no Referencial Curricular do Paraná, efetivam o prescrito na LDB/96/97 e são propositivas ao indicar a importância de : 1.
Conhecer os aspectos estruturantes das diferentes tradições/organizações \textbf{religiosos} e filosofias de vida, a partir de pressupostos científicos, filosóficos, estéticos e éticos. 2.
Compreender, valorizar e respeitar as manifestações \textbf{religiosas} e filosofias de vida, suas experiências e saberes, em diferentes tempos, espaços e \textbf{territórios}. 3.
Reconhecer e cuidar de si, do outro, da coletividade e da natureza, enquanto expressão de valor da vida. 4.
Conviver com a diversidade de crenças, pensamentos, convicções, modos de ser e viver. 5.
Analisar as relações entre as tradições \textbf{religiosas} e os campos da cultura, da política, da economia, da saúde, da ciência, da tecnologia e do meio ambiente. 6.
Debater, problematizar e posicionar -se frente aos discursos e práticas de \textbf{intolerância}, \textbf{discriminação} e violência de cunho \textbf{religioso}, de modo a assegurar os direitos humanos no constante exercício da cidadania e da cultura de paz. ( BNCC, BRASIL. 2017, pg. 435 ).
Dessa forma, as Competências Gerais e Específicas propostas para o Ensino \textbf{Religioso} foram contempladas e tratadas no âmbito dos Direitos e Objetivos de aprendizagem.
Por conseguinte, as Unidades Temáticas correlacionam -se entre si e recebem ênfases diferentes, de acordo com cada ano de escolarização.
Os Objetos de Conhecimento são os conhecimentos básicos essenciais que os estudantes têm direito de aprender e que são desdobrados em Objetivos de Aprendizagem.
Assim, tendo em vista a trajetória do \textbf{Estado} do Paraná e de alguns de seus Municípios no que diz respeito à experiência com o componente Ensino \textbf{Religioso}, na proposta do presente documento se inserem Objetos de Conhecimento complementares, relacionados com a Unidade Temática, a fim de favorecer a transição dos Anos Iniciais para os Anos Finais do Ensino Fundamental, e, também, por uma abordagem \textbf{hierarquizada} de objetos de conhecimento, ampliando gradativamente o nível de aprendizagem.
Procurou -se superar a fragmentação dos conhecimentos e a ruptura dos mesmos na transição do Ensino Fundamental - Anos Iniciais e Finais, sendo proposto para cada ano, um conjunto progressivo de conhecimentos historicamente construídos, de forma que o estudante tenha um percurso contínuo de aprendizagem.
Nessa perspectiva, os objetos de conhecimento foram ampliados em praticamente todos os anos, permitindo que o processo de aprendizagem e desenvolvimento da educação no Ensino Fundamental possam ser contempladas integralmente.
As Unidades Temáticas que compõem a BNCC e, portanto, constam no Referencial Curricular do Paraná são : \textbf{Identidades} e alteridades ; Manifestações \textbf{religiosas} ; Crenças \textbf{Religiosas} e Filosofias de Vida.
A partir dessas Unidades Temáticas, foram estabelecidos na BNCC, os objetos de conhecimento para cada ano, que são : práticas espirituais ou ritualísticas, espaços e \textbf{territórios} sagrados, mitos, crenças, narrativas, oralidade, tradições orais e textos escritos, doutrinas, ideias de imortalidade ( \textbf{ancestralidade}, reencarnação, ressurreição, transmigração, entre outras ), códigos éticos e filosofias de vida.
Sendo assim, os critérios de organização das habilidades na BNCC ( com a explicitação dos objetos de conhecimento aos quais se relacionam e do agrupamento desses objetos em Unidades Temáticas ) expressam um arranjo possível, dentre muitos outros, para a realidade de cada \textbf{Estado} e Município da Federação.
Ao considerar as especificidades da disciplina, ressalta -se que os encaminhamentos metodológicos devem primar pela garantia dos direitos de aprendizagem e estar em consonância com a legislação vigente.
Ademais, a avaliação deve ser concebida sob uma perspectiva formativa com a finalidade de acompanhamento do processo de ensino e aprendizagem.
Ressalta -se que, para o desenvolvimento do encaminhamento pedagógico em sala de aula, os professores contemplem as quatro matrizes que formam a religiosidade brasileira : Matriz \textbf{Indígena}, Matriz Africana, Matriz Ocidental e Matriz Oriental.
O estudo destas matrizes tem por objetivo fortalecer o exercício da cidadania, o fomento ao conhecimento, além de ampliar os horizontes dos estudantes em relação à diversidade \textbf{religiosa}.
O diálogo inter-religioso é uma possibilidade de superação do grande desafio da humanidade : vivermos juntos em paz com respeito e alteridade. [ 24 ] : Conforme o Dicionário de Filosofia Nicola Abbagnano a etimologia da palavra \textbf{laico} tem origem no termo Grego laon ( adj : laikós - λαϊκό ), expressão que designava o povo em sentido lato, tão abrangente ou tão universal quanto possível.
O termo laon, ou laikós referia -se, portanto, à entidade do cidadão, \textbf{população}, ao povo todo, a toda a gente, sem exceção alguma.
Também pode ser encontrado em dicionários como sendo o “ \textbf{laico} ” uma forma erudita e “ leigo ” a forma vulgar ; ambas vieram do Latim LAICUS, do Grego LAIKOS, LAICISMO ( in Laicism- ingl, fr.
Laicisme ; it.
Laicismó ). ( ABBAGNANO, 2003 ).

% matched lemmas: ancestralidade, discriminação, estado, hierarquizar, identidade, indígena, intolerância, laico, população, religioso, território
\end{textsample}
